\documentclass[11pt]{article}
\usepackage{fontspec}
\usepackage{amsmath,amssymb,amsthm,mathtools}
\usepackage[margin=1.1in]{geometry}
\usepackage{microtype}
\usepackage{parskip}
\usepackage{lmodern}
\usepackage[colorlinks=true,linkcolor=blue!50!black,urlcolor=blue!50!black]{hyperref}
\providecommand{\tightlist}{}
\title{Quantum theory from sealed records II: fermions, statistics, and internal symmetry groups as outputs}
\author{Felix Robles Elvira (ORCID: 0009-0009-2017-4394; independent researcher)}
\date{\small preprint, not peer reviewed, version 2026-06-11.}
\begin{document}
\maketitle

\section*{Abstract}
Within the sealed-records program (Part I), we derive the structures of particle statistics and gauge theory as outputs rather than inputs.  Results at stated scope, each with machine receipts: (i) \textbf{univalence superselection is derived within the stated record class} (defined in Section 2; the silence of the relative sign is established by exhaustive verification within that class, so the verb "derived" carries exactly that scope) — the silent-seam axiom then excludes superpositions across the double cover, making the $(-1)^F$ grading central; (ii) \textbf{spin quantization} $2m \in \mathbb Z$ follows from the $(1,3)$ signature with the belt-trick identification, with record anyons confined to genuinely two-dimensional screens; (iii) the \textbf{statistics theorem} — \emph{under the named framed-transport premise} (exchange is transport of the fiber with its frame; graded as premise throughout, with a corpus-bound [P] derivation not part of this batch), sealed positivity on the exchange collar forces exchange phase $(-1)^{2m}$: wrong-statistics sectors carry strictly negative weights, and Pauli exclusion appears as an exact zero; (iv) \textbf{parastatistics excluded at pair-exchange level}: the transport commutant on $V \otimes V$ is exactly $\mathrm{span}\{\mathbf 1, P\}$, and the $\pm\mathbf 1$ options are excluded by the ledger (the $n$-strand DHR parastatistics question is scoped in Section 4); (v) the \textbf{internal symmetry group of the fiber is reconstructed as the commutant of record exchange}: once the ledger forces the statistics, classical Schur–Weyl duality (named import, no novelty claimed) delivers the full unitary group of the fiber — $U(3)$ for three-strand fibers (implementation check: $165 = 165$ dimension match, determinant slot included; the baryon channel transforms by $\det U$ exactly), $U(2)$ for the weak fiber ($10 = 10$, $20 = 20$) — in the \emph{direction} of Doplicher–Roberts: group as output of statistics.  What is delivered is the \emph{global} symmetry group; the promotion to a local gauge symmetry with connection and dynamics is an additional step, graded as input where Part III uses it.  Color receipts include exact Pauli for color ($\Vert A_4\Vert = 0$) and the uniqueness of the baryon singlet.  We state the paper's two named premises — framed transport, carrying the exchange association, and collar presentability, named in Section 3's premise accounting — and the scope of each result; the matter-content selection built on these structures, and its falsifiable consequences, appear in Part III.

\section*{1. Introduction}
Standard quantum field theory \emph{postulates} its deepest structural facts: that particles are bosons or fermions with the spin-statistics connection, that parastatistics does not occur, and that interactions are governed by compact gauge groups acting on internal indices.  The algebraic tradition showed these postulates are not independent — Doplicher–Haag–Roberts theory [1] reconstructs a compact gauge group from the category of superselection sectors — but the inputs there are themselves field-theoretic axioms.

This paper derives the same structures inside the sealed-records program (Part I): the only primitives are records and the three axioms R, S, C.  The derivations are \emph{exclusion} arguments in the program's characteristic style — the silent-seam axiom S eliminates every alternative — and each step carries machine receipts.

\section*{2. Univalence superselection, derived}
\textbf{Theorem 2.1 (stated scope).}  Half-integer record fibers are consistent (the silent-sign analysis: a $2\pi$ frame rotation's sign is unobservable on any single record), but \emph{superpositions across the double cover are excluded}: a relative sign between integer and half-integer sectors would be a silent seam — a datum no record registers — and axiom S quotients it away.  Hence the grading $(-1)^F$ is superselected: central in the record algebra.  The argument form is Wick–Wightman–Wigner's univalence superselection [8], and we credit it as such; the program's specific contribution is the conversion of "unobservable" into "excluded by axiom S" via the operational criterion below, with its negative branch verified within the stated record class (the abstract's claim carries that scope).

\textbf{The stated record class (defined, since the criterion below quantifies over it).}  Throughout this paper the record class is the finite sealed-record towers of Part I Section 2: finite products of projective record fibers with sealed seams, closed under the dilation lifts of this section, probed exhaustively up to the stated tower depths in the receipts.  Claims of silence are exhaustive-verification statements \emph{within this class}, graded as such — silence in a finite class is evidence for, not a proof of, silence in every extension, and we say so once here instead of hedging every sentence.

\textbf{The silent-vs-physical criterion (stated, since it carries the load here and in Part I).}  A relative datum between two ledger configurations is \emph{physical} if some admissible sealed record assigns them different weights, and \emph{silent} if no record does; axiom S quotients silent data away.  The criterion is operational, not aesthetic: each application below is discharged by a receipt exhibiting either the discriminating record (physical) or the exhaustive failure to find one within the stated record class (silent).  The $2\pi$ rotation sign is silent on every single record (receipted); a relative phase \emph{within} a sector changes seam interference weights (receipted), so intra-sector superpositions survive; the integer/half-integer relative sign fails the test against every record — that asymmetry is the theorem.

The receipts include the $4\pi$-interferometry consistency checks and the projective-dilation classification ($H^2$ classes, both realized).  Two consequences do heavy work later: (a) fermionic intertwiners cannot implement sector equivalences (they shift a superselected label) — load-bearing in Part III's charge-screening theorem; (b) spin quantization: with the $(1,3)$ signature of Part I and exchange-as-transport (below), $2m \in \mathbb Z$, with anyonic interpolation surviving only on genuinely two-dimensional record screens.

\section*{3. The statistics theorem}
\textbf{The association.}  The first of this paper's two named premises (the second, collar presentability, is named in the premise accounting below): exchange of identical record structures is \emph{framed transport} — the exchange path is the transport of one fiber around the other with its frame.  (In the program's corpus this premise is discharged into a no-silent-circulation theorem for eventless transport on open paths: silent extra circulation is quantized to $\mathbb Z_2$ and a silent $\pi$ is excluded.  That derivation is \textbf{corpus-bound [P]} — not in this batch — so within this submission the framed-transport association is a \emph{named premise}, full stop; the [P] pointer records where a derivation exists to be audited, and nothing here leans on it.)

\textbf{Theorem 3.1 (with its premises displayed).}  Sealed positivity on the exchange collar forces the exchange phase $(-1)^{2m}$. Wrong-statistics sectors are not merely unconventional: their collar weights are strictly negative — the same failure class as the reflection-positivity violations of Part I — and Pauli exclusion appears as an exact zero of the sealed weight.

\emph{Premise accounting, complete.}  (a) The framed-transport association above.  (b) \textbf{The presentation premise, previously silent and now named}: Part I proves reflection positivity as a theorem only for eventless sectors admitting finite primitive Markov presentations, and proves that hypothesis sharp; applying "sealed positivity" on the exchange collar therefore assumes the collar sector \emph{admits such a presentation}.  Our collar constructions do (they are built as finite record lattices, and the receipts certify positivity directly on them), but the general statement "every physical exchange collar is presentable" is a premise of the same type as the tame premise of Part I.  (c) We also temper the novelty claim: "wrong statistics implies positivity violation" is the engine of the classical field-theoretic proof [3]; what differs here is the input layer (records, not fields) and the finite, machine-checkable form of the violation.

Receipts (implementation checks in the sense of Section 5, except the first, which is the substantive one): the wrong-sector negativity is exhibited directly; the fermionic tower then reconstructs canonical anticommutation relations to $8.9\times 10^{-16}$ with Fock subset-sum structure at $7.6\times 10^{-15}$, and a quadratic reconstruction whose inputs nowhere assume the CAR algebra (that is what "non-circular" means here) at $2.4\times 10^{-15}$.

\section*{4. Parastatistics never, and why dimension three matters}
\textbf{Theorem 4.1.}  (a) For relative configuration spaces of dimension $\ge 3$, eventless exchange squares to the identity (pure-gauge holonomy receipts at $2.2\times 10^{-14}$ across contraction stages; an eventful contrast degrades from $1.99$ to $0.12$ — holonomy is \emph{enclosed events}).  The configuration-space formulation of exchange used throughout this section is Leinaas–Myrheim's [5] (named frame, not only the source of the $d = 2$ exception).  (b) The transport commutant on $V \otimes V$ is exactly $\mathrm{span}\{\mathbf 1, P\}$ (machine dimension $2$ at $d = 2, 3, 4$), so the exchange operator is $\pm\mathbf 1$ or $\pm P$; the scalar options are \emph{excluded by the ledger}: they predict vanishing cross-records at seams where eventless transport in fact seals with unit modulus — a distinction between labelings that no record registers, which is precisely the silent-datum configuration the criterion of Section 2 excludes.  Hence exchange $= (-1)^{2m}P$ at pair level: ordinary Bose/Fermi statistics, with the two-dimensional anyon window [5] as the unique exception.

\emph{Scope of the parastatistics exclusion, stated precisely.}  The theorem above is a statement about \emph{pair exchange}.  Parastatistics in the DHR sense lives in higher-dimensional representations of $S_n$ acting on multiplicity spaces, and excluding it requires the $n$-strand statement.  What this batch has: (i) the pair-level theorem; (ii) the $n = 3, 4$ channel receipts of Section 5 (the statistics algebra truncates as ordinary Bose/Fermi structure requires — rank $23 \neq 24$ on $(\mathbb C^3)^{\otimes 4}$ — and sector bookkeeping matches $3^n$ counting), which are \emph{consistency evidence at the probed strand numbers, not a proof}; and (iii) the observation that a genuine para-order-2 alternative is excluded at $n = 3$ by the baryon-channel computation of Section 5.  The all-$n$ exclusion theorem is open and named; the abstract is scoped accordingly.

\section*{5. The gauge group as the commutant of exchange}
\textbf{Theorem 5.1 (gauge commutant; Schur–Weyl applied, with the physical identification named).}  The algebra of record-transport operators commuting with all exchanges on $n$ strands of a $d$-dimensional fiber is exactly $\mathrm{span}\{U^{\otimes n} : U \in U(d)\}$ — the \emph{full} unitary group of the fiber, determinant slot included.

\emph{Status of the statement.}  The mathematics here is classical Schur–Weyl duality [4]; we prove nothing new about commutants and claim no mathematical novelty.  What the program contributes is the two-step physical identification that makes the classical theorem do gauge-theoretic work: (i) the ledger \emph{forces} the statistics (Sections 3–4), so the symmetric-group action is derived, not assumed; (ii) the named premise that internal symmetries are record-transport operators, which must therefore commute with the forced statistics — whereupon Schur–Weyl delivers the group.  The analogy with Doplicher–Roberts reconstruction [1, 2] is one of direction (group as output), not of mathematical content or depth.

Receipts: for the color fiber ($d = 3$, $n = 3$), commutant dimension $165$ equals the dimension of $\mathrm{span}\{U^{\otimes 3}\}$ ($165 = 10^2 + 8^2 + 1^2$); the statistics algebra truncates exactly as Schur–Weyl requires (rank $23 \neq 24$ on $(\mathbb C^3)^{\otimes 4}$); for the weak fiber ($d = 2$), the mirror receipts give $10 = 10$ ($n = 2$) and $20 = 20$ ($n = 3$).  \emph{These numbers verify our implementation against a known theorem — they are consistency checks of the machinery, not evidence for new mathematics}; the receipts that carry evidential weight in this section are the ledger-side ones (negativity of wrong-statistics sectors, the baryon-channel weight $0.0096 > 0$ below), where the program's claims are not classical. Gauge structure is thereby an \emph{output}: the group is whatever commutes with the statistics the ledger forces.

\textbf{Color receipts.}  The three-strand antisymmetric channel is unique ($\operatorname{rank} A_3 = 1$: one baryon singlet); the four-strand channel vanishes \emph{exactly} ($\Vert A_4\Vert = 0$: Pauli for color); and the baryon transforms by $\det U$ alone ($U^{\otimes 3} b = \det(U)\, b$ to $2.4\times 10^{-16}$) — the determinant slot that Part III identifies as the hypercharge carrier.

\textbf{Order parameters of the derivation.}  Fiber order equals fiber dimension, checked in both directions \emph{within the constructed model}: a parastatistics-order-2 alternative predicts zero sealed weight in the baryon channel, where the constructed record dynamics yields $0.0096 > 0$ — a model output demonstrating that the record construction realizes the dimension-equals-order case, not a measurement of nature.  Sector collapse bookkeeping matches $3^n$-counting exactly (same status).

\section*{6. What is not claimed}
(i) The specific fiber dimensions $\{3, 2\}$ are \emph{selected}, not derived ab initio: the program's exhaustion theorem makes $3$ the first chirality-supporting stack and a tower-exclusion search makes $(3,2)$ the strictly cheapest chiral tower on uniformly derived charge lattices — with the minimality principle itself named as the remaining input (Part III).  (ii) Dynamics (running couplings, mass generation) is outside this paper.  (iii) All results are at stated scope, with the program's frontier taxonomy (rank-free vs engine-uniform vs frontier-bound) applying clause by clause; notably the centrality and commutant arguments here are \emph{rank-free} — their proofs never invoke the fiber multiplicity.

\section*{7. Relation to prior work}
DHR/DR reconstruction [1, 2] is the conceptual ancestor; the difference is the input layer (records and exclusions rather than field-theoretic axioms) and the direction: statistics is \emph{forced} (Section 3) before the group is reconstructed from it (Section 5). The spin-statistics literature is engaged at two levels: the field-theoretic theorem [3], and the long program of seeking kinematic or non-relativistic derivations, whose history and pitfalls are documented by Duck and Sudarshan [6] and whose best known geometric attempt is Berry–Robbins [7] — our no-silent-circulation route belongs to this second family and inherits its burden: the exchange-as-framed-transport association is a named premise, not a theorem of quantum mechanics.  The two-dimensional anyon window is Leinaas–Myrheim and Wilczek [5]. Schur–Weyl duality [4] supplies the double-commutant machinery (named import, with the honesty note of Section 5).

\section*{Reproducibility}
Receipts regenerate from fixed-seed scripts (the silent-sign and holonomy analyses; the commutant dimensions; the CAR tower; the baryon channel).  Bit-identical reruns; repository accompanies the program.

\section*{References}
[1] S. Doplicher, R. Haag, J. E. Roberts, Local observables and particle statistics I–II, \emph{Commun. Math. Phys.} \textbf{23} (1971) 199, \textbf{35} (1974) 49.

[2] S. Doplicher, J. E. Roberts, Why there is a field algebra with a compact gauge group describing the superselection structure in particle physics, \emph{Commun. Math. Phys.} \textbf{131} (1990) 51.

[3] R. F. Streater, A. S. Wightman, \emph{PCT, Spin and Statistics, and All That}, Princeton (1964).

[4] R. Goodman, N. Wallach, \emph{Symmetry, Representations, and Invariants}, Springer (2009) — Schur–Weyl duality.

[5] J. M. Leinaas, J. Myrheim, On the theory of identical particles, \emph{Nuovo Cim. B} \textbf{37} (1977) 1; F. Wilczek, Quantum mechanics of fractional-spin particles, \emph{Phys. Rev. Lett.} \textbf{49} (1982) 957.

[6] I. Duck, E. C. G. Sudarshan, \emph{Pauli and the Spin-Statistics Theorem}, World Scientific (1997).

[7] M. V. Berry, J. M. Robbins, Indistinguishability for quantum particles: spin, statistics and the geometric phase, \emph{Proc. R. Soc. Lond. A} \textbf{453} (1997) 1771.

[8] G. C. Wick, A. S. Wightman, E. P. Wigner, The intrinsic parity of elementary particles, \emph{Phys. Rev.} \textbf{88} (1952) 101 — univalence superselection.

\textbf{Part I:} \emph{Quantum theory from sealed records I} (same author). \textbf{Part III:} \emph{The Standard-Model floor} (same author).

\end{document}
